\section{Discussion}
\label{sec:discussion}

The convergence theory of Gaussian natural gradient flows separates into three
stages, and the separation persists from the continuous dynamics to the
deterministic discretizations and to the stochastic algorithms. An underdispersed
initial covariance costs a burn-in time logarithmic in
$\beta \lambda_{0,\min}$; the global phase localizes the iterate at the intrinsic
condition number obtained after whitening by the optimizing covariance; and the
local phase runs at a rate set by the spectral gap of a self-adjoint score
operator rather than by the curvature ratio of the target. The three stages have
different sharpness certificates. The logarithmic spiral shows that the
$\kappastar$ dependence of the global phase cannot be improved under the general
spectral initialization, the bump train shows that the factor $\kappa/\Delta t$
of the fixed-step schemes is sharp for a constant-factor reduction, and the
convex ridge shows that the $\betastar-1$ branch of the local spectral scale is
attained, so that no dimension-free logarithmic local rate can hold.

The stochastic theory preserves this structure with the sticking-the-landing and
Price/Hessian estimators, at the cost of explicit variance floors: the
Hessian-Lipschitz floors scale as $1/B$ in the batch size, while the
assumption-minimal global floor plateaus at $O(\kappa N_\theta)$ for
$B\lesssim\kappa$ and the local floor at $O(N_\theta^2\LHs^2/V_1)$ for
$B\lesssim V_1$; all of these floors vanish for Gaussian targets. The stochastic local result is a killed-process contraction together with a
finite-horizon exit-probability bound; because the mean-score noise has
unbounded support, we do not claim pathwise invariance of a bounded local region
at a fixed stepsize, and we do not claim infinite-horizon containment. The
discrete local radius for a general target is constructive through a supremum of
second derivatives of the update map, and is not a closed-form quantity that a
practitioner can evaluate.

Several questions remain open. The local spectral bound degrades with dimension
through the factor $\sqrt{N_\theta} \log \kappastar$, and the convex-ridge
family lives in dimension
$N_\theta = \Theta(\kappastar^2)$. We do not know whether a stronger bound holds
in fixed dimension, nor whether the $\sqrt{N_\theta} \log \kappastar$ branch is
itself sharp. The fixed-step lower bound is stated for one globally
fixed stepsize, and it leaves open whether an adaptive, line-search, or implicit
scheme can evade the factor $\kappa/\Delta t$. A closed-form discrete local
radius, and sharper stepsize windows for the stochastic local phase, would make
the local theory usable as a stopping rule rather than as an existence
statement. We prove a stationarity guarantee for the covariance-constrained KL
scheme in the deterministic setting only; its stochastic counterpart requires a
separate treatment of the interaction between the clipping step and the noise,
which we do not treat. Finally, the classification of
\Cref{sec:geometry} exhibits a three-parameter family of affine-invariant
metrics whose modal rates we compute exactly on Gaussian targets, but the sharp
global and local theory of \Cref{sec:global,sec:local} is developed for
Fisher--Rao alone; extending it to unbalanced members of the family, and
developing the Wasserstein--Fisher--Rao dynamics, which the transport warm start
of \Cref{sec:global} does not develop, are natural directions for subsequent
work.
